\scanpage{24}
% archive sheet number: 23

est orth ssi elle satisfait à $\lambda\rho'+1=0$ (cas $n$ pair) resp $\underbrace{\lambda\rho+2}_{=\lambda\rho'+2}=0$ (cas $n$ impair).

Donc

a) Si $n$ pair, $\exists$ une telle $u$ \inserted{$\ne\mathrm{id}$ sur tte fibre} ssi $\rho'$ inv, et on a alors unicité :
\[
u=\mathrm{id}-\pi\otimes o'\qquad o'=\frac1{\rho'}\widetilde o\tag*{(6)}
\]
i.e. $u$ est l'\uline{antipodisme}

b) Si $n$ impair, \inserted{$2$ inv}, \deleted{$\exists$ une telle $u$ ssi} \inserted{pour qu'il} \deleted{$\rho$ est \unsure{inv|inv.}, $\exists!u$} \deleted{ou bien $\rho$.} existe une telle $u$, il \deleted{faut} \inserted{f. et} suffit que $\rho' (=\rho^{\deleted{\unsure{\prime|*}}})$ soit inv, \deleted{et $2$ inv.} \deleted{cette condition est néc.} \deleted{et alors cette condition est} \deleted{\unsure{néc. si} $2$ inv.} On a alors unicité,

% margin note written sideways
\margin{dans ce cas, $u$ \unsure{invarie|inverse|invariant} donc automatiquement la \unsure{réal.|virt.} \unsure{géom|géom.} envisagée de $\Pi$}

$\lambda=-\dfrac2\rho=-\dfrac2{\rho'}$ i.e.
\[
u=\mathrm{id}-\pi\otimes\frac{2\widetilde o}{\rho'}\tag*{(7)}
\]
et $u$ est encore l'\uline{antipodisme}.

c) Si \inserted{$n$ impair,} $S$ de car. $2$, \deleted{alors $\exists$ un} \deleted{ssi $\rho$ est} \inserted{pour qu'il existe,} une telle $u$, il \deleted{\unsure{faut}} faut que $\rho(=\rho')$ \deleted{soit} \unsure{nilpotent}, et suffit que \deleted{$\rho=\rho'$ soit \unsure{nul}} \deleted{\unsure{soit inv.}} Dans ce dernier cas, \unsure{tte} pseudo reflexion qui est identité sur $V$ et commute à $G$ est orthogonale,

% slanted
\margin{(La condition néc et suff. est donc $\rho'=0$, \deleted{\unsure{dans}} \inserted{\unsure{le cas|les cas}} \unsure{lorsque} $S$ réduit)}

\deleted{en particulier}

En particulier

\uline{Corollaire} Supposons $\sigma'\rho'$ inv. i.e. $f(\widetilde o)$ inv. i.e. $c\ \deleted{\unsure{\notin}}\in\Gamma(X\setminus Q_X)$, \deleted{en ttes fibres} i.e. $o\notin Q_X$ en tte fibre, \inserted{\uline{et} $2$ inv. sur $S$ si $n$ impair.} Alors, et alors seulement, \deleted{lors que $G$} $\exists$ unique $G$-automorphisme $\ne\mathrm{id}$ sur chaque fibre, induisant l'identité sur $C$. et invariant $Q_C$. Ce dernier est l'antipodisme.
